%----------------------------------------------------------------------------------------------
%-----------------------------------     4. Integralrechnung     ------------------------------
%----------------------------------------------------------------------------------------------
\section{Integralrechnung}{493}

Eine Funktion $F(x)$ heißt Stammfunktion von $f(x)$, wenn ihre Ableitung wieder $f(x)$ ist. Das unbestimmte Integral von $f(x)$ ist dann

$$ \int f(x) \,  \diff x = F(x) + C  \quad \quad \Longleftrightarrow \quad \quad F'(x) = \frac{\diff \:F(x)}{\diff x} = f(x)$$

Beim Ableiten der Stammfunktion verschwindet die Konstante $C$, da die Ableitung einer Konstanten $0$ ist.


%----------------------------------------------------------------------------------------------
%--------------------------------     4.1 Obersumme / Untersumme     --------------------------
%----------------------------------------------------------------------------------------------
\subsection{Obersumme / Untersumme}

Die Fläche unter einer Funktion $f(x)$ wird in Intervalle zerlegt.

\textbf{Voraussetzung: $f: [a;b] \to \mathbb{R}$ und $\mathbb{W}_f$ beschränkt} 

\renewcommand{\arraystretch}{1.5}
\begin{tabular}{@{}ll@{}}
    Zerlegung                           & $Z = \lbrace x_0; x_1; \ldots ; x_n \rbrace = \lbrace x_i$  $\vert i=0; \ldots ; n \rbrace$ \quad $(n \in \mathbb{N}; n \geq 2)$ \\
    Breite(n) der Intervalle            & $\Delta x_i = x_{i+1} - x_i$ \quad $\Delta x$ kann variieren! \\
    Maschenweite                        & $\mathrm{max}(\Delta x_i) = d(Z)$ \\
    \textbf{Äquidistante Zerlegung}     & $Z = \lbrace a + \frac{b-a}{n} \cdot i $ $\vert i=0; \ldots; n \rbrace$
\end{tabular}

\vspace{0.4cm}

\begin{minipage}[c]{0.48\columnwidth}
    \textbf{Untersumme \myul{US}} 
    $$ \rm US = \sum \limits_{i=0}^{n-1} m_i \cdot \Delta x_i $$
    $$ m_i = \mathrm{inf} \lbrace f(x) \, \vert x_i \leq x \leq x_{i+1} \rbrace $$

    \begin{center}
        \includegraphics[width=0.42\columnwidth]{images/untersumme.jpg}
    \end{center}
\end{minipage}
\hfill
\begin{minipage}[c]{0.48\columnwidth}
    \textbf{Obersumme $\overline{\rm OS}$} 
    $$ \rm OS = \sum \limits_{i=0}^{n-1} M_i \cdot \Delta x_i $$
    $$ M_i = \mathrm{sup} \lbrace f(x) \, \vert \,  x_i \leq x \leq x_{i+1} \rbrace $$

    \begin{center}
        \includegraphics[width=0.42\columnwidth]{images/obersumme.jpg}
    \end{center}
\end{minipage}


\begin{minipage}{0.45\columnwidth}
    \textbf{Unter-Integral}

    \renewcommand{\arraystretch}{2.5}
    \begin{tabular}{l}
        $I = \lim \limits_{\diff (Z) \to 0} US $ \quad \textcolor{blue}{$\infty \cdot 0$} \\
        $I = \int \limits_a^{b} f(x) \, \diff x $
    \end{tabular}
    \renewcommand{\arraystretch}{1}
\end{minipage}
\hfill				
\begin{minipage}{0.45\columnwidth}
    \textbf{Ober-Integral}

    \renewcommand{\arraystretch}{2.5}
    \begin{tabular}{ll}
        $I = \lim \limits_{\diff (Z) \to 0} OS$ \quad \textcolor{blue}{$\infty \cdot 0$} \\
        $I = \int \limits_{a}^{b} f(x) \, \diff x $ 
    \end{tabular}
    \renewcommand{\arraystretch}{1}
\end{minipage}


%----------------------------------------------------------------------------------------------
%---------------------------------     4.2 /3  Riemann-Summe     ------------------------------
%----------------------------------------------------------------------------------------------
\subsection{Riemann-Summe}{506-507}

\vspace{-0.2cm}
$$ \text{RS} = \sum \limits_{i=0}^{n-1} f(\xi) \, \Delta x_i  \qquad \diff (Z) \to 0 \qquad \int \limits_{a}^{b} f(x) \,  \diff x \qquad\xi \in [x_i ; x_{i+1}] $$	

\subsection{Riemann-Integral}{507}

\begin{minipage}{0.4\columnwidth}
    $$ \text{Notation:} \quad \int \limits_{a}^{b} f(x) \, \diff x $$
\end{minipage}				
\hfill
\begin{minipage}{0.55\columnwidth}
    \begin{tabular}{ll}
                & beschränkt $[a;b]$\\
    Bedingung:  & bedingt Stetig\\
                & wenige unstetigkeit stellen
    \end{tabular}
\end{minipage}	


%----------------------------------------------------------------------------------------------
%-----------------------------     4.4 Wichtige Integrale S. 495     --------------------------
%----------------------------------------------------------------------------------------------
\subsection{Wichtige Integrale S. 495}{495}

\begin{minipage}[c]{0.26\columnwidth}
    $$ \int \limits_{a}^{b} x^2 \, \diff x = \frac{b^3}{3} - \frac{a^3}{3} $$
\end{minipage}
\hfill
\begin{minipage}[c]{0.26\columnwidth}
    $$ \int \limits_{a}^{b} x \, \diff x = \frac{b^2}{2} - \frac{a^2}{2} $$
\end{minipage}
\hfill
\begin{minipage}[c]{0.35\columnwidth}
    $$ \int \limits_{a}^{b} 1 \, \diff x = b - a  \text{ (Rechteck)} $$ 
\end{minipage}


%----------------------------------------------------------------------------------------------
%-----------------------------     4.5 Flächen unter Integralen     ---------------------------
%----------------------------------------------------------------------------------------------
\subsection{Flächen unter Integralen}

\textbf{Voraussetzung: $f: [a;b] \to \mathbb{R}$ und $\mathbb{W}_f$ beschränkt}
$$ \text{Fläche:} \quad A = \int \limits_{a}^{b} \abs{ f(x) } \, \diff x $$
\textcolor{red}{Der Inhalt des Betrags muss auf Vorzeichen untersucht werden!} \\
Negative Vorzeichen müssen über die x-Achse geklappt werden: \\
$\vert x^2 - x \vert = -(x^2 - x)$ falls $x > x^2$ 


%----------------------------------------------------------------------------------------------
%----------------------------------     4.6 Integralregeln     --------------------------------
%----------------------------------------------------------------------------------------------
\subsection{Integralregeln}{496-499}

\subsubsection{Elementare Regeln An1b}

\vspace{-0.2cm}
\renewcommand{\arraystretch}{1.6}
\begin{tabular}{llll}
    Konstanten:     & $f(x) = c \quad$ konst                & $\int f(x) \diff x$                      & $F(x) = cx + C$\\
    Faktor:         & $f(x) = \alpha \cdot x^\beta$             & $\alpha \int f(x) \diff x$           & $F(x) = \alpha \cdot \frac{x^\beta}{\beta} + C$ \\
    Summen:         & $f(x) = u(x) + v(x) $                 & $\int f(x) \diff x$                      & $\int u(x) \diff x + \int v(x) \diff x$    
\end{tabular}


%----------------------------------------------------------------------------------------------
%------------------------------     4.7 Wichtige Integrale    ---------------------------------
%----------------------------------------------------------------------------------------------
\subsection{Wichtige Integrale}{446}

\begin{tabular}{|ll|ll||ll|ll|}
    \hline
    Funktionen                                 && Integrierung                             && Funktionen                                 && Integrierung                                                    &\\
    \hline
    $ x^n$                                     && $ \frac{x^{n+1}}{n+1} $                  && $ \frac{1}{cosh^2 x} $                     && $ \tanh x $                                                     &\\
    $ \frac{1}{x} $                            && $ \ln \abs{x} $                          && $ \frac{1}{\sinh^2 x}$                     && $ -\coth x $                                                    &\\
    $ e^x $                                    && $ e^x $                                  &&                                            &&                                                                 &\\
    $ a^x $                                    && $ \frac{a^x}{\ln a} $                    && $ \frac{1}{\alpha^2 + x^2}$                && $ \frac{1}{\alpha} \arctan \frac{x}{\alpha} $                   &\\
    $ \sin x $                                 && $ -\cos x $                              && $ \frac{1}{\alpha^2 - x^2}$                && $ \frac{1}{2\alpha} \ln \frac{\alpha + x}{\alpha - x} $         &\\
    $ \cos $                                   && $ \sin x $                               &&                                            && $ (\abs{x} < \alpha, \; 0 < \alpha)$                            &\\
    $ \tan x $                                 && $ -\ln \abs{\cos x} $                    && $ \frac{1}{x^2 - \alpha^2}$                && $ \frac{1}{2\alpha} \ln \frac{\alpha - x}{\alpha + x} $         &\\
    $ \cot x $                                 && $ \ln \abs{\sin x}$                      &&                                            && $ (\abs{x} > \alpha, \; 0 > \alpha)$                            &\\
    $ \frac{1}{cos^2 x} $                      && $ \tan x $                               && $ \frac{1}{\sqrt{\alpha^2 + x^2}}$         && $ \ln \abs{x + \sqrt{\alpha^2 + x^2}}$                          &\\
    $ \frac{1}{\sin^2 x}$                      && $ -\cot x $                              &&                                            && $ (0 < a)$                                                      &\\
    $ \sinh x$                                 && $ \cosh x $                              && $ \frac{1}{\sqrt{\alpha^2 - x^2}}$         && $ \arcsin \frac{x}{\alpha} $                                    &\\
    $ \cosh x$                                 && $ \sinh x $                              &&                                            && $ (\abs{x} < \alpha, \: 0 < \alpha) $                           &\\
    $ \tanh x $                                && $ \ln \abs{\cosh x} $                    && $ \frac{1}{\sqrt{x^2 - \alpha^2}}$         && $ \ln \abs{x + \sqrt{\alpha^2 - x^2}} $                         &\\
    $ \coth x $                                && $ \ln \abs{\sinh x} $                    &&                                            && $ (\abs{x} > \alpha, \: 0 > \alpha) $                           &\\
    \hline
\end{tabular}


%----------------------------------------------------------------------------------------------
%--------------------------------     4.8 Mittelwertsatz    -----------------------------------
%----------------------------------------------------------------------------------------------	
\subsection{Mittelwertsatz}{510}

\textbf{Voraussetzung: $f: [a;b] \to \mathbb{R}$ und $\mathbb{W}_f$ beschränkt und stetig}

Die Fläche unter der Funktion $f(x)$ kann an mind. einem Punkt $\xi$ als Rechteck dargestellt werden. \\
Es muss gelten: $\xi \in (a;b)$

\begin{minipage}[c]{0.40\columnwidth}
    \includegraphics[width=0.95\columnwidth]{images/MWS-integral.PNG}
\end{minipage}
\hfill				
\begin{minipage}[c]{0.55\columnwidth}
    \raggedright
    \textbf{Fläche = Mittelwert $\bm{\cdot}$ Intervall-Länge}
    
    $\frac{1}{b-a} \int \limits_{a}^{b} f(x) \, dx = f(\xi)$ 

    $ \int \limits_{a}^{b} f(x) \, \diff x = (b-a) \cdot f(\xi)$
    $ (b-a) \cdot h = (b-a) \cdot f(\xi)$ 
\end{minipage}
